
% ===========================================================================
% Abstract
% ===========================================================================
\begin{center}
\textbf{Abstract:}
\end{center}
\begin{quote}
\small
Blockchain-based reputation systems for additive manufacturing (AM) supply
chains must withstand strategic adversarial manipulation: actors who
misreport quality, forge identities, collude to amplify scores, or exploit
access-control weaknesses to gain unearned trust.  Formal game-theoretic
analysis of such multi-agent strategic interactions is intractable for
realistic system parameters, yet prior evaluations rely exclusively on
static test scripts that probe a fixed set of known attack patterns.  This
paper presents a complementary methodology in which \emph{learned
adversaries} trained simultaneously with honest agents using Multi-Agent
Proximal Policy Optimization (MAPPO) continuously search for exploitable
weaknesses across a 12-action, 14-dimensional observation space.  We model
a Bayesian Beta-reputation system with Wilson 95\% confidence intervals,
write-time exponential temporal decay, and stake-based economic incentives,
and evaluate it across 11 experimental configurations spanning 9 distinct
attack classes: self-rating manipulation, Sybil flooding, collusion rings,
administrative privilege escalation, evidence tampering, reputation-gate
bypass, provenance replay, and combinations thereof.  Trained across 5
independent random seeds per configuration (10{,}000 episodes per seed),
the honest population converges to $\geq\!99.7\%$ honest behavior in 8 of 9
single-attack configurations.  Deterministic defense mechanisms achieve
100\% block rates for self-rating, administrative escalation, and gate bypass
attacks; probabilistic detection ($p=0.80$) yields an 81.5\% empirical block
rate for evidence tampering; and economic deterrence constrains Sybil and
collusion exploitation to 52--58\% without eliminating those strategic modes.
The comprehensive multi-vector configuration identifies the primary open
challenge: concurrent multi-vector attacks degrade honest convergence to
$69.9\%$ in the training tail with a 54.1\% defense hold rate across
134{,}337 attack attempts.  All source code, configurations, and training
logs are released as a reproducible reference implementation.
\end{quote}

\vspace{4pt}
{\small\textbf{Keywords:} multi-agent reinforcement learning; adversarial
evaluation; blockchain; reputation system; Bayesian inference; additive
manufacturing; supply chain security; MAPPO; PettingZoo; defense mechanisms}
\vspace{10pt}
\hrule

% ===========================================================================
% I. Introduction
% ===========================================================================
\section{Introduction}
\label{sec:intro}

Additive manufacturing (AM) supply chains are inherently multi-stakeholder:
material certification labs, print service bureaus, post-processing vendors,
quality inspectors, logistics providers, and OEMs each contribute to and
depend upon the collective trustworthiness of the supply chain.
Blockchain-based reputation systems have been proposed as a mechanism to
encode this distributed trust into a tamper-evident ledger, so that future
supply chain participants can condition their cooperation decisions on the
demonstrated performance history of their peers~\cite{almaayn2024unified,
nakamoto2008bitcoin, androulaki2018hyperledger}.

A reputation system is only as valuable as its resistance to manipulation.
An actor who can inflate their own score through self-rating, or who can
suppress competitors' scores through collusion or Sybil attacks, undermines
the very signal the system is designed to provide.  The canonical approach
to adversarial evaluation is scripted attack injection: a tester writes
fixed exploit sequences targeting known vulnerabilities, executes them
against a running system, and reports which defenses hold~\cite{ferretti2021blockchain,
douceur2002sybil}.  This methodology is rigorous for the attacks it covers
but cannot discover emergent strategies that the analyst did not anticipate.

\emph{Multi-Agent Reinforcement Learning} (MARL) offers a complementary
evaluation paradigm~\cite{lowe2017multi, foerster2018counterfactual,
yu2021surprising}.  When adversarial agents learn their policies by trial
and error, competing against a simultaneously learning honest population,
they effectively conduct a continuous automated search for exploitable
weaknesses.  Defenses that are robust under MARL training provide stronger
robustness guarantees than defenses that hold only against hand-crafted
exploits.

\subsection{Problem Statement}

We address the following research question: \emph{Can a Bayesian
Beta-reputation system with stake-based economic incentives, Wilson
confidence-interval scoring, and layered access controls withstand
adversarial agents that learn optimal manipulation strategies through
multi-agent reinforcement learning?}  Specifically:
\begin{enumerate}[label=\textbf{RQ\arabic*.}, leftmargin=*, noitemsep]
  \item Do honest agents converge to high honest-action rates despite the
        presence of adversarial peers, and how does convergence depend on the
        type and intensity of attack?
  \item Which defense mechanisms provide deterministic versus probabilistic
        protection, and what residual exploitation rate do remaining attack
        vectors achieve after convergence?
  \item At what adversarial population fraction does multi-vector attack
        coordination overwhelm the defense layer, and does the system degrade
        gracefully or catastrophically?
\end{enumerate}

\subsection{Contributions}

This paper makes the following contributions:
\begin{itemize}[noitemsep]
  \item A \emph{PettingZoo AEC environment} modeling a Bayesian reputation
        system with 12 discrete actions encompassing all 9 attack vectors
        from the companion blockchain evaluation~\cite{almaayn2024unified}
        and a 14-dimensional observation space (\S\ref{sec:env}).
  \item A \emph{MAPPO training harness} with shared policy networks, GAE
        advantage estimation, and per-seed convergence tracking, enabling
        reproducible multi-seed evaluation across 11 experimental
        configurations (\S\ref{sec:marl}).
  \item \emph{Empirical defense characterization} across three classes:
        deterministic penalty ($-5.0$ per attempt), probabilistic detection
        ($p=0.80$), and economic deterrence; quantified as block rates over
        10{,}000-episode runs with 5 independent seeds each (\S\ref{sec:eval}).
  \item Identification of the \emph{comprehensive multi-vector attack} as the
        primary open challenge, with honest convergence degrading to $69.9\%$
        and 54.1\% defense hold rate across 134{,}337 attack attempts
        (\S\ref{sec:discussion}).
\end{itemize}

\subsection{Paper Organization}

Section~\ref{sec:related} reviews related work.
Section~\ref{sec:sysmodel} describes the reputation system model.
Section~\ref{sec:marl} details the MARL framework and training procedure.
Section~\ref{sec:eval} presents experimental results.
Section~\ref{sec:discussion} discusses key findings and limitations.
Section~\ref{sec:conclusion} concludes.
